% TOP_PROBLEM: {"problemId": "problem.two-dimensional-jacobian-conjecture-in-characteristic-zero", "canonicalTitle": "Two-Dimensional Jacobian Conjecture in Characteristic Zero", "releaseRank": 103, "primaryDomain": "algebra_representation_category", "primaryDomainLabel": "Algebra, representation & category theory", "displayStatus": "open", "releaseStatus": "open"}
% !TeX program = lualatex
% Definitions and sources retained from research_batch_101_103.tex; research is in GPT_6_Astra_Ultra.
\documentclass[11pt]{article}
\usepackage[margin=25mm]{geometry}
\usepackage{fontspec}
\setmainfont{Latin Modern Roman}
\usepackage{amsmath,amssymb,mathtools}
\usepackage{unicode-math}
\setmathfont{Latin Modern Math}
\usepackage{enumitem,needspace}
\usepackage{xurl,hyperref}
\hypersetup{colorlinks=true,urlcolor=blue}
\setcounter{secnumdepth}{1}
\setlength{\parindent}{0pt}
\setlength{\parskip}{5pt}
\setlength{\emergencystretch}{3em}
\setlist[itemize]{leftmargin=3em,itemsep=5pt,topsep=4pt,parsep=0pt}
\allowdisplaybreaks
\newcommand{\N}{\mathbb{N}}
\newcommand{\Z}{\mathbb{Z}}
\newcommand{\Q}{\mathbb{Q}}
\newcommand{\R}{\mathbb{R}}
\newcommand{\C}{\mathbb{C}}
\newcommand{\F}{\mathbb{F}}
\newcommand{\A}{\mathbb{A}}
\newcommand{\PP}{\mathbb P}
\newcommand{\E}{\mathbb{E}}
\newcommand{\eps}{\varepsilon}
\newcommand{\dd}{\,\mathrm d}
\newcommand{\sref}[2]{\hyperlink{source:#1:#2}{[S#2]}}
\newcommand{\eref}[2]{\hyperlink{extra:#1:#2}{[E#2]}}
% Turn the next switch off for a shorter reading copy. The quotations preserve
% every exactTarget field, including qualifications omitted from a short summary.
\newif\ifshowcatalogue
\showcataloguetrue
\newcommand{\cataloguescope}[1]{\ifshowcatalogue
 \paragraph{Exact scope in the supplied catalogue.}
 \begin{quote}\small #1\end{quote}\fi}
\begin{document}
\setcounter{section}{102}
% ================================================================
% problemId: problem.two-dimensional-jacobian-conjecture-in-characteristic-zero
% source releaseRank: 103; source releaseStatus: open
% exactTarget SHA-256: 8ea39d9d3dc7447f107db63684a8c632dfb14985b3201d5472b02cc82b029d57
\clearpage
\section{\texorpdfstring{Two-Dimensional Jacobian Conjecture in Characteristic Zero}{Two-Dimensional Jacobian Conjecture in Characteristic Zero}}\label{103}
\subsection{Definitions and mathematical statement}
For $F=(f,g)$ with $f,g\in k[x,y]$, its Jacobian determinant is
\[
 JF=\frac{\partial f}{\partial x}\frac{\partial g}{\partial y}
       -\frac{\partial f}{\partial y}\frac{\partial g}{\partial x}.
\]
A Keller map satisfies $JF\in k^\times$. In characteristic zero, the question asks whether there are $u,v\in k[x,y]$ making $G=(u,v)$ a two-sided polynomial inverse: $G\circ F=F\circ G=\operatorname{id}$. This is the affine plane case over every characteristic-zero field. Statements or alleged counterexamples in more variables do not decide this selected two-variable target.
\par\smallskip\textit{Formulation and concept references:} \sref{103}{1}.
\cataloguescope{Let k be a field of characteristic zero and let F:k\textasciicircum{}2 -> k\textasciicircum{}2 be a polynomial map whose Jacobian determinant is a nonzero constant. Determine whether F must have a polynomial inverse, equivalently whether every two-variable Keller map is a polynomial automorphism of the affine plane.}
\subsection{Short English statement}
Does every plane polynomial map with constant nonzero Jacobian determinant have a polynomial inverse?
\subsection{Sources}
\begin{itemize}\raggedright
\item[\textnormal{[S1]}]\hypertarget{source:103:1}{} Terence Tao, "A digestion of the Jacobian conjecture counterexample", What's new, 21 July 2026.
\url{https://terrytao.wordpress.com/2026/07/21/a-digestion-of-the-jacobian-conjecture-counterexample/}.
{\small\textit{Catalogue formulation source.}}
\item[\textnormal{[S2]}]\hypertarget{source:103:2}{} Shuhong Gao, Counterexamples to the Jacobian conjecture in dimensions greater than two.
\url{https://arxiv.org/abs/2608.00222}.
{\small\textit{Catalogue research/status reference; not a proof endorsement.}}
\item[\textnormal{[S3]}]\hypertarget{source:103:3}{} Gao, section 2: Background and previous work.
\url{https://arxiv.org/html/2608.00222v1#S2}.
{\small\textit{Catalogue research/status reference; not a proof endorsement.}}
\item[\textnormal{[E1]}]\hypertarget{extra:103:1}{}
J. A. Guccione, J. J. Guccione, R. Horruitiner, and C. Valqui,
\emph{Increasing the degree of a possible counterexample to the Jacobian
Conjecture from 100 to 108}, arXiv:2204.14178v1, 29 April 2022.
Theorem 2.1 and Sections 2--6; the total-degree exclusion is distinct from
this notebook's degree-in-one-variable calculation.
\url{https://arxiv.org/html/2204.14178v1}.
{\small\textit{Consulted 25 September 2026; cited as a partial-result
manuscript, not as a proof of the full conjecture.}}
% END RESEARCH ADDITION REFERENCES-103
\end{itemize}

\end{document}
